\documentclass[10pt,a4paper]{article}
\usepackage[margin=1.82cm]{geometry}
\usepackage[round,authoryear]{natbib}
\usepackage{times}
\usepackage[T1]{fontenc}
\usepackage{textcomp}
\usepackage{microtype}
\usepackage[hidelinks]{hyperref}
\setlength{\parindent}{1.05em}
\setlength{\parskip}{0pt}
\pagestyle{plain}

\begin{document}

\begin{center}
{\Large\bfseries Executive Summary}\\[0.45em]
{\large\bfseries  Periodic Timetable Modelling \& 
Matheuristic Event-Retiming Optimization\\[0.3em]
of the Swiss Passenger Rail Network}\\[0.45em]
Matthias Benjam\'in Andrews Garc\'ia
\end{center}

Swiss rail planning is shaped by infrastructure, rolling stock, capacity projects, stations, and long-term programmes such as Bahn 2000 and STEP 2035. The passenger value of these investments, however, is ultimately experienced through the timetable. In a clockface network, the timetable forms a repeated time-space structure that determines which services can be reached, how much time is spent on board, how much is absorbed by same-train dwell and transfers, and how long passengers wait before a useful departure. Studying the timetable therefore provides a direct way to measure how infrastructure and operating choices translate into passenger accessibility. This perspective follows the accessibility literature, which evaluates transport systems through the opportunities and generalized costs they create rather than through infrastructure extent alone \citep{hansen1959accessibility,geurs2004accessibility}. It is particularly relevant in Switzerland, where recurring services and coordinated nodes make passenger benefits sensitive to small timing relationships.

The thesis has two linked parts. Part I is descriptive: it reconstructs historical Swiss rail timetables and asks how passenger-weighted journey conditions have evolved over time and across space. Part II is prospective: it tests whether small integer-minute retimings can reduce passenger-minutes in the 2026 and 2035 timetables while preserving section running times and the recognizable structure of the existing timetable. The two parts use the same conceptual unit, namely a routed passenger journey decomposed into rolling time, same-train dwell, transfer time, initial waiting, and total post-departure time. The first part therefore measures where passengers lose time; the second part searches for localized timetable changes that can recover some of that time.

Part I combines compact periodic timetable-token expansion, schedule-based routing, OD weighting, generalized-cost projection, and iterative proportional fitting. A 2023 NPVM public-transport OD matrix is used as a demand anchor, then converted to station-to-station weights and projected across historical years \citep{vmuvek2026npvm2023}. The resulting OD matrices are weighting layers for comparing timetable states. The matrix construction follows the usual separation between a structured spatial-interaction seed and external marginal constraints \citep{wilson1970entropy,demingStephan1940IPF,fienberg1970IPF}. The routing layer retains nondominated paths and allocates demand by dominance minutes, meaning that paths covering a larger portion of the clockface departure interval receive a larger share of the OD flow. This approach follows schedule-based public-transport assignment, where actual departures, transfers, and attractive alternatives matter for route choice \citep{spiessFlorian1989OptimalStrategies,nuzzoloRussoCrisalli2001ScheduleAssignment}.

The main historical result is a substantial but uneven improvement. For the selected 131-station network, passenger-weighted post-departure time falls from 51.60 minutes per trip in 1982 to 39.13 minutes in 2026, while initial waiting falls from 28.51 to 13.53 minutes. These gains reflect rolling-time reductions, improved frequencies, shorter transfers, changed dwell patterns, and better dominance of attractive departures. Several timetable windows allow for a temporal analysis. The Bahn 2000 period improved average travel time the most; the base-tunnel periods produced strong axis-specific gains; and the 2021--2026 window improved initial waiting while worsening post-departure travel times on average, as a result of 2025's punctuality-oriented timetable relaxation of Western Switzerland \citep{sbbRomandie2025Punctuality}. National means therefore hide both spatial redistribution and component trade-offs.

Part II expands the model from 131 selected stations to nearly the full Swiss station universe: 1'592 stations in 2026 and 1'579 in 2035. The timetable inputs combine the public network diagrams and the STEP 2035 planning horizon, while the demand layer uses SBB station-frequency observations where available, modelled fallbacks elsewhere, a gravity-shaped OD seed, a transfer-throughput correction, and IPF balancing \citep{smaPartnerDownloads,bav2021as2035,sbbPassagierfrequenzDataset2026}. Within each model year, demand is held fixed across the Baseline, Optimized Step 1, Optimized Step 2, and Optimized Step 3 scenarios, because the goal is to isolate timetable effects. Every completed scenario is evaluated by rebuilding retained nondominated OD paths and reallocating the fixed OD demand through the dominance-minute rule.

The Part II validation objective is the total daily passenger-minute cost of the modelled Swiss OD system after departure. For each origin–destination pair, fixed daily demand is multiplied by the dominance-weighted average cost of the retained nondominated paths under a given timetable. Path cost includes rolling time, same-train dwell, and transfer time, while pre-departure initial waiting time is reported separately. A timetable edit is therefore valuable only if it reduces the weighted sum of routed journey costs after passengers are reallocated across the paths made attractive or unattractive by the timetable. Local retimings are judged by their effect on the whole routed OD system: predetermined OD demand levels, path choice, and subsequent network-wide rerouting all determine whether a connection-level improvement becomes a systemwide passenger-minute saving.

The optimization problem is deliberately narrower than a full national timetable redesign. Compared to the other selected literature, the approach of this thesis is a local-search event-retiming matheuristic. It formulates a restricted 0--1 integer-programming candidate-selection problem, then implements it through a heuristic local-search procedure in which candidates are evaluated by a schedule-based passenger-routing and travel-time model. Mathematical-programming approaches such as PESP and Benders decomposition can provide strong formulations and, for suitable instances, optimality bounds \citep{serafini1989periodic,liebchen2007pesp,leutwilerCorman2022LogicBenders}. Passenger-oriented decomposition and column-generation methods show how routing can be integrated more directly \citep{schiewe2020integrated,martinIradiRopke2022Matheuristic,yaoEtAl2025HybridBenders}. Trepat Borecka et al. further show that even the choice of geographic decomposition boundaries is itself an important algorithmic decision \citep{trepatBoreckaEtAl2025Geographic}. The thesis applies a model-based matheuristic because it is most compatible with the scale of the national station coverage, roughly 2.5 million OD pairs per model year, dominance-based path allocation, and full-network validation. Matheuristics are suitable here because mathematical objectives, constraints, and exact evaluations can guide a heuristic search without requiring the complete national problem to be solved globally, leading to substantial results within a reasonable runtime \citep{boschettiManiezzo2022Matheuristics}.

The search proceeds through three stages. Step 1 screens high-flow existing transfers and near misses using station-board proxy scores. Step 2 evaluates additional hub-filtered candidates using exact paths retained under the validated Step 1 timetable. Step 3 broadens the candidate pool, lowers the transfer-context threshold, repeatedly re-scores remaining candidates after each accepted edit, and reviews the final combined result. The workflow separates iterative candidate selection from final scenario evaluation. Step 3 candidate selection took 8 hours for 2026 and 5.2 hours for 2035 on an Apple M1 Max with 64 GB RAM, while a complete two-year routing rebuild of the final Step 3 timetable required 133.8 core-hours, or 14.1 hours using ten parallel processes, before demand enrichment and analysis of the OD matrices.

The final corrected Step 3 timetable improves the national passenger-time objective in both years. In 2026, 42 edited lines save 94'216 passenger-minutes per day, equivalent to 2.78 seconds across all trips, or 73.9 seconds across only adjusted OD pairs. In 2035, 37 edited lines save 112'419 passenger-minutes per day, equivalent to 3.12 seconds per weighted trip, or 83.0 seconds across only adjusted OD pairs. Using year-corrected Swiss public-transport values of travel-time savings \citep{axhausenEtAl2004SwissVTTS}, these savings correspond to approximately CHF 7.6 million per year in 2026 and CHF 13.3 million per year in 2035. Most savings arise from transfer and same-train dwell reductions. Rolling time can rise in aggregate because full OD rerouting reallocates passengers to different retained paths; this does not mean trains run more slowly, since section running times are preserved by construction. It means passengers would prefer a geographical diversion if a transfer or dwell reduction along that route lowers total post-departure time.

The distributional result is as important as the national mean. Only 3.76\% of positive-demand OD pairs change in 2026 and 3.38\% in 2035. The improvement is therefore selective rather than uniform. Among OD pairs with at least one passenger per day, 2026 contains 1'635 OD pairs with improved existing transfers averaging -1.10 minutes per trip and 736 OD pairs with worsened existing transfers averaging +0.74 minutes. OD pairs defined by newly unlocked transfers are fewer, 211, but average -1.76 minutes. Newly missed transfers account for only 111 OD pairs, averaging +1.04 minutes. In 2035, there are 1'873 OD pairs with improved existing transfers, averaging -1.72 minutes, while 431 OD pairs with newly unlocked transfers average -5.26 minutes; newly worsened transfers number only 107 and average +0.96 minutes. OD pairs with changed retained-path counts require a separate interpretation: additional alternatives can appear because a new path becomes attractive, but the dominance-weighted mean can still worsen if shares shift toward a slower retained option. Conversely, fewer retained paths can improve the OD mean when an inferior or less useful option disappears from the nondominated set. The value of a retiming therefore depends on the sign, size, demand, and mechanism of each changed relation.

Spatially, the effects are concentrated and highly varied. Ticino has the strongest 2026 origin-canton improvement, while Solothurn leads in 2035. Station and municipality effects are more varied still. Reachability and passenger-minute rankings also measure different outcomes: in 2035, Schlieren reaches 1'192 destinations one minute faster, yet saves only about one quarter of Olten's saved passenger-minutes because Olten's 101 affected destinations carry greater demand. Segment-load analysis shows fixed demand shifting between IC and S-Bahn alternatives in Zurich, between S10 and S20 in Ticino, and between IC61/IC82 and S26/S29. These results show that local retimings can affect path choice and passenger loads far beyond the edited sections, even though most OD pairs remain unchanged.

The methodological contribution is an integrated national application rather than a new universal exact optimization algorithm. The thesis connects accessibility measurement, OD construction, schedule-based routing, dominance allocation, periodic-event retiming, and matheuristic candidate selection into one workflow. It shows that the well-defined historical timetable travel-time measurements can become an intervention framework: the same path and component representation that is used to explain long-run changes can identify transfer-sensitive structures and evaluate candidate retimings. The practical contribution is a compact shortlist of 42 line changes in 2026 and 37 line changes in 2035 that have passed structural validation, sequential compatibility checks, and full passenger-network rerouting. The accepted line edits are each linked to a passenger-minute score, compatibility context, and downstream OD effect.

The limitations are important to consider: the optimized timetables are passenger-weighted timetable hypotheses instead of operational implementation plans, meaning that platform occupation, microscopic headways and capacity constraints, rolling-stock circulation, crew duties, stochastic delay propagation, fares, and induced demand remain outside the model. A railway operator would still need to verify whether each line edit is feasible under local infrastructure, and resource constraints. However, because most timing shifts are merely one or two minutes, it is likely that a large majority of the edits would pass such checks. The thesis also holds demand fixed within each model year, which is appropriate for isolating timing effects but not for predicting induced demand. Future work can therefore proceed in two complementary directions: operational validation of the accepted edits, and more advanced decomposition or incremental-routing methods that reduce the cost of evaluating passenger responses. Within these boundaries, the core conclusion is clear: even an already highly coordinated Swiss clockface timetable contains measurable passenger-time potential in small transfer-sensitive retimings, and that potential is made visible when timetable edits are evaluated through the improved passenger paths they facilitate.

\clearpage
\bibliographystyle{plainnat}
\bibliography{executive_summary}

\end{document}
